\chapter{Marco práctico: ¿aporta valor supervisar a un agente LLM?}
\label{cap:marco-practico}

% Fuente única de cifras: notebooks/STRATA_marco_practico.ipynb y los JSON que carga
% (decision_automl_prep.json, panel_mm25_inclGBM-XGB-SE_AUC_emb1_N0-150_step21_kfold_seed42.json,
% automl_importance.json, strategy_clustering15.json, spy_ablation_robustness.json,
% m10_smci_valtest_robustez.json). Toda cifra de este capítulo está trazada a esos ficheros.
% Brief del capítulo: memoria/MARCO_PRACTICO_SPEC.md. Figuras: por exportar desde el notebook.

Este trabajo pregunta si una capa de supervisión estadística clásica logra rescatar a un agente de trading cuando falla, y si ese rescate puede probarse. El agente que estudiamos, cinco personalidades sobre un modelo de lenguaje, decide cada día una dirección y un tamaño de posición. Sobre el periodo de evaluación pierde dinero y acierta la dirección menos de la mitad de las veces. STRATA observa cada decisión con tres detectores de series temporales, régimen, cambio de opinión y volatilidad, y produce señales que tres estrategias derivadas explotan: la regla a mano M8, el meta-learner XGBoost M10 y una búsqueda automática de modelo con AutoML.

El capítulo sostiene que supervisar al agente añade valor medible, y lo contrasta por partes. Las estrategias supervisadas baten al agente en acierto direccional con significancia y le recortan el riesgo. Un aprendiz potente, libre de elegir entre muchas familias de modelos, reencuentra las señales de STRATA antes que inventar otras. Y al agrupar los activos por su naturaleza asoma qué estrategia conviene a cada tipo. El límite va por delante: batir a las estrategias triviales en accuracy es nominal con esta muestra. STRATA no promete alfa; rescata al perdedor y acota dónde funciona.

\section{Datos, universo y protocolo}
\label{sec:mp-datos}

El caso central es \textbf{SPY}, el ETF del S\&P 500, acompañado de un panel de robustez de quince activos que abarca índices amplios, sectoriales, acciones individuales, un proxy de criptomoneda y una materia prima. La elección de SPY como escaparate descansa en el \emph{leverage effect} (Black 1976 \cite{black1976}; Christie 1982 \cite{christie1982}): en un índice amplio la volatilidad alta coincide con las caídas \emph{de forma contemporánea}, de modo que el régimen de volatilidad del HMM adquiere contenido de riesgo sobre el que la supervisión puede actuar. Conviene acotar el alcance desde el principio: esa coincidencia es del mismo día y no convierte al régimen en un predictor de la dirección del día siguiente. La asunción se debilita además en valores individuales con apalancamiento flojo, y esa frontera la documentamos más adelante.

Los tres detectores se calibran una sola vez sobre el periodo 2000-01-01 a 2024-09-30 (veinticuatro años) y sus umbrales quedan fijados \emph{ex ante}, nunca sobre los datos de evaluación (Sección~\ref{sec:strata-calib}). La evaluación es fuera de muestra desde 2024-10-01 hasta el cierre del trabajo, con inicio posterior al corte de conocimiento del modelo de lenguaje que mueve al agente, lo que descarta que esas decisiones se memorizaran en su entrenamiento. Tras el arranque del walk-forward quedan en torno a doscientos cincuenta días de evaluación por activo.

El protocolo desplegable es un walk-forward expandible con arranque de ciento cincuenta días, reentreno cada veintiuno y embargo de un día, justificado por el horizonte de la etiqueta en un esquema de origen rodante (Tashman 2000 \cite{tashman2000}; López de Prado 2018 \cite{lopezdeprado2018}). La convención causal es la de la Sección~\ref{sec:cpcv-causalidad}: la posición del día $t$ multiplica al retorno del día $t+1$, con \texttt{signal\_lag}${}=1$.

\section{Mecánica ex ante: los tres detectores}
\label{sec:mp-mecanica}

STRATA combina tres detectores ortogonales, todos calibrados antes de mirar la evaluación. RAM lee el régimen con un HMM gaussiano de tres estados (Hamilton 1989 \cite{hamilton1989}) sobre el log-retorno y la volatilidad realizada a veintiún días, y mide si la dirección del agente concuerda con el régimen; su prior régimen-a-signo se calibra por activo a partir de las medias del periodo de ajuste. GSO aplica un GARCH(1,1) de innovaciones Student-t (Bollerslev 1986 \cite{bollerslev1986}) y juzga si el tamaño es compatible con la volatilidad prevista. PSA, más sutil, detecta con el algoritmo bayesiano de Adams y MacKay cuándo el agente cambia de opinión de forma estructural. Las definiciones formales y los umbrales de los tres están en las Secciones~\ref{sec:strata-ram}, \ref{sec:strata-gso} y~\ref{sec:strata-psa}.

La Figura~\ref{fig:mp-regimen-spy} muestra el régimen filtrado del HMM sobre la curva de SPY en el periodo de evaluación: los tramos de Crisis se concentran en las caídas, manifestación contemporánea del \emph{leverage effect}. La precisión importa para no prometer de más. En la calibración el retorno del mismo día baja de forma monótona con el régimen ($+0{,}00054$ en Calma, $+0{,}00017$ en Estrés, $-0{,}000002$ en Crisis), pero el retorno del día siguiente es positivo en los tres regímenes y la fracción de días al alza ronda el azar (en torno a $0{,}52$ en Crisis): el régimen separa por volatilidad y coincide con las caídas el mismo día, no anticipa el signo de mañana. Su valor para la supervisión está en disciplinar el riesgo, no en predecir la dirección.

% [Figura por exportar — notebook §4.2, celda 6: curva B&H de SPY (OOS) coloreada por régimen
%  filtrado del HMM (Calma/Estrés/Crisis). Fuente: build_states("SPY").]
\begin{figure}[htbp]
\centering
\fbox{\parbox{0.92\textwidth}{\centering\vspace{1.2cm}\textit{Régimen HMM filtrado sobre la curva de SPY (OOS), exportar desde el notebook \texttt{STRATA\_marco\_practico.ipynb} §4.2.}\vspace{1.2cm}}}
\caption{Régimen del HMM (Calma, Estrés, Crisis) sobre el nivel de SPY en el periodo de evaluación. Las fases de Crisis coinciden con las caídas, condición del \emph{leverage effect}.}
\label{fig:mp-regimen-spy}
\end{figure}

\section{Caso de estudio: SPY}
\label{sec:mp-spy}

SPY es el escaparate del mecanismo. El agente solo, que denotamos M5, acierta el $36{,}6\%$ de las direcciones, muy por debajo del azar, con un Sharpe de $-3{,}07$ y un capital final del $0{,}70$ del inicial. El sign test contra $0{,}5$ rechaza la hipótesis de no-habilidad con $p < 0{,}001$: el agente es un perdedor direccional claro, que es el punto de partida del problema.

La Tabla~\ref{tab:mp-spy} recoge las seis estrategias sobre SPY. La regla M8 sube la accuracy a $0{,}442$ y recorta el daño en riesgo (Sharpe $-0{,}46$, máxima caída del $15{,}2\%$ frente al $30{,}2\%$ del agente). El meta-learner M10 alcanza $0{,}494$. La búsqueda automática AutoML llega a $0{,}574$ y gana en punto a todas las estrategias, incluidas las triviales, con un Sharpe de $+2{,}68$. Las triviales comprar-y-mantener y clase mayoritaria coinciden aquí en $0{,}566$, porque SPY fue alcista en la ventana y la clase dominante es estar largo.

\begin{table}[htbp]
\centering
\caption{SPY, periodo de evaluación walk-forward (config.\ mm25, embargo${}=1$). Accuracy direccional, Sharpe, máxima caída, Calmar y capital final. Fuente: \texttt{panel\_mm25\_\ldots\_seed42.json}.}
\label{tab:mp-spy}
\begin{tabular}{lrrrrr}
\hline
Estrategia & Accuracy & Sharpe & Máx.\ caída & Calmar & Capital \\
\hline
\textbf{AutoML (búsqueda H2O)} & \textbf{0,574} & \textbf{+2,68} & $-5{,}5\%$ & 6,93 & 1,38$\times$ \\
ZeroR (clase mayoritaria) & 0,566 & +2,21 & $-9{,}8\%$ & 3,10 & 1,30$\times$ \\
B\&H (comprar y mantener) & 0,566 & +2,21 & $-9{,}8\%$ & 3,10 & 1,30$\times$ \\
M10 (meta-learner XGBoost) & 0,494 & $-0{,}60$ & $-16{,}1\%$ & $-0{,}50$ & 0,92$\times$ \\
M8 (regla STRATA) & 0,442 & $-0{,}46$ & $-15{,}2\%$ & $-0{,}39$ & 0,94$\times$ \\
M5 (agente solo) & 0,366 & $-3{,}07$ & $-30{,}2\%$ & $-1{,}00$ & 0,70$\times$ \\
\hline
\end{tabular}
\end{table}

Ese ``gana a todo'' de AutoML admite dos lecturas, y las dos son ciertas. El McNemar pareado sobre los aciertos diarios da $p = 0{,}007$ para M10 frente a M5 y $p = 0{,}0002$ para AutoML frente a M5: STRATA rescata al agente con significancia. El mismo test de AutoML frente a ZeroR da $p = 0{,}90$, sin potencia para separarlos. La ventaja en accuracy sobre el mercado pasivo es nominal, fruto de una ventana de evaluación corta, y así la reportamos.

El rescate también se ve en el riesgo. Las curvas de capital de la Figura~\ref{fig:mp-equity-spy} muestran al agente hundiéndose mientras M8 y M10 controlan la caída. A nivel de un solo activo el intervalo de confianza del bootstrap estacionario (Politis y Romano 1994 \cite{politis1994}) sobre la diferencia de Sharpe todavía cruza el cero, porque doscientos cincuenta días dan poca potencia; la significancia del rescate de riesgo aparece al agrupar el panel (Sección~\ref{sec:mp-panel}).

% [Figura por exportar — notebook §4.3, celda 12: curvas de capital de M5, M8, M10, ZeroR y B&H sobre SPY.]
\begin{figure}[htbp]
\centering
\fbox{\parbox{0.92\textwidth}{\centering\vspace{1.0cm}\textit{Curvas de capital (M5, M8, M10, ZeroR, B\&H) sobre SPY, exportar desde el notebook §4.3.}\vspace{1.0cm}}}
\caption{Capital acumulado de las estrategias sobre SPY en el periodo de evaluación. El agente solo se hunde; M8 y M10 contienen la caída.}
\label{fig:mp-equity-spy}
\end{figure}

La importancia por valores de Shapley sobre el mejor árbol del leaderboard de AutoML y la importancia por permutación del ensemble coinciden en colocar las señales de STRATA por delante de las del agente, primer indicio de la universalidad que el panel confirma.

\section{Generalización y universalidad (panel de quince activos)}
\label{sec:mp-panel}

Para distinguir el efecto de la suerte llevamos el análisis a los quince activos. La Tabla~\ref{tab:mp-panel} reúne la accuracy de cada estrategia por activo, y su lectura conjunta es matizada. En los índices amplios (SPY, QQQ, DIA, IWM, XLK) las triviales suelen ganar en accuracy, porque fueron mercados alcistas y la clase mayoritaria cuesta de batir. Donde el activo es volátil o de apalancamiento débil, en cambio, asoma una derivada de STRATA que supera a las dos triviales a la vez: M8 en MSTR, M10 en SMCI y en UNG, AutoML en MARA. El rescate del agente es, con todo, transversal: alguna derivada de STRATA mejora a M5 en casi todo el panel.

\begin{table}[htbp]
\centering
\caption{Accuracy direccional por activo y estrategia (panel de quince, walk-forward, embargo${}=1$). En negrita, por activo, la mejor derivada de STRATA cuando supera a las dos triviales. El número de días de evaluación varía por activo (246 a 259) según la fecha de arranque del walk-forward. Fuente: \texttt{panel\_mm25\_\ldots\_seed42.json}.}
\label{tab:mp-panel}
\small
\begin{tabular}{lrrrrrr}
\hline
Activo & M5 & M8 & M10 & AutoML & ZeroR & B\&H \\
\hline
SPY  & 0,366 & 0,442 & 0,494 & 0,574 & 0,566 & 0,566 \\
QQQ  & 0,418 & 0,486 & 0,522 & 0,534 & 0,590 & 0,590 \\
DIA  & 0,440 & 0,484 & 0,468 & 0,520 & 0,552 & 0,556 \\
IWM  & 0,450 & 0,470 & 0,458 & 0,482 & 0,554 & 0,574 \\
XLE  & 0,448 & 0,528 & 0,508 & 0,532 & 0,565 & 0,565 \\
XLF  & 0,429 & 0,502 & 0,526 & 0,510 & 0,538 & 0,538 \\
XLK  & 0,510 & 0,470 & 0,542 & 0,590 & 0,641 & 0,641 \\
NVDA & 0,467 & 0,521 & 0,483 & 0,517 & 0,552 & 0,552 \\
BAC  & 0,451 & 0,516 & 0,419 & 0,488 & 0,561 & 0,561 \\
TSLA & 0,474 & 0,478 & 0,522 & 0,454 & 0,458 & 0,530 \\
MSTR & 0,554 & \textbf{0,558} & 0,534 & 0,498 & 0,530 & 0,450 \\
SMCI & 0,484 & 0,496 & \textbf{0,552} & 0,472 & 0,516 & 0,484 \\
ROKU & 0,444 & 0,528 & 0,508 & 0,544 & 0,548 & 0,548 \\
MARA & 0,528 & 0,528 & 0,532 & \textbf{0,544} & 0,532 & 0,468 \\
UNG  & 0,510 & 0,502 & \textbf{0,518} & 0,482 & 0,449 & 0,486 \\
\hline
\end{tabular}
\end{table}

La aportación dura del trabajo está en el riesgo agregado. Agrupando los retornos diarios de los quince activos (tres mil setecientos cincuenta y un días) y aplicando un bootstrap estacionario pareado por bloques (Politis y Romano 1994 \cite{politis1994}), el rescate de riesgo de la regla frente al agente es significativo: M8 mejora el Sharpe de M5 en $+0{,}66$ con intervalo de confianza al $95\%$ de $[0{,}23,\ 1{,}16]$, y reduce la máxima caída en $+0{,}24$ con intervalo $[0{,}02,\ 0{,}44]$; ambos excluyen el cero (Tabla~\ref{tab:mp-pooled}). El mismo contraste para M10 frente a M5 queda en el límite, con el intervalo del Sharpe rozando el cero. La supervisión disciplina el riesgo del agente con evidencia que sobrevive a un test.

\begin{table}[htbp]
\centering
\caption{Rescate de riesgo agregado (bootstrap estacionario pareado por bloques, $n=3751$ días, quince activos). Diferencias frente al agente M5 con intervalo al $95\%$. Fuente: \texttt{decision\_automl\_prep.json}, bloque \texttt{pooled}.}
\label{tab:mp-pooled}
\begin{tabular}{lrrc}
\hline
Comparación & $\Delta$Sharpe & IC$_{95}$ & ¿excluye 0? \\
\hline
M8 vs M5  & $+0{,}66$ & $[0{,}23,\ 1{,}16]$ & sí \\
M10 vs M5 & $+0{,}61$ & $[-0{,}04,\ 1{,}25]$ & no \\
\hline
 & $\Delta$Máx.\ caída & IC$_{95}$ & ¿excluye 0? \\
\hline
M8 vs M5  & $+0{,}24$ & $[0{,}02,\ 0{,}44]$ & sí \\
M10 vs M5 & $+0{,}23$ & $[-0{,}02,\ 0{,}45]$ & no \\
\hline
\end{tabular}
\end{table}

La universalidad se mide con la importancia por características y con una ablación. Repartiendo la importancia por valores de Shapley (Lundberg y Lee 2017 \cite{lundberg2017}) en bloques, las características de STRATA concentran de media el $66\%$ del total (régimen $0{,}38$, volatilidad $0{,}15$ y cambio de opinión $0{,}14$, frente a $0{,}34$ del agente), y la cuota de STRATA supera la mitad en los quince activos. El aprendiz potente, lejos de descartar las señales de STRATA, se apoya en ellas. La ablación añade un matiz honesto: sumar las siete señales de STRATA al vector de quince características del agente no eleva la accuracy del meta-learner de media ($-0{,}05$ en SPY, $+0{,}08$ en SMCI). El valor de STRATA está en el rescate y en la interpretabilidad.

% [Figura por exportar — notebook §4.4, celdas 17-18: SHAP medio por bloque (MEDIA DEL PANEL, no SPY)
%  y heatmap de accuracy por activo y estrategia centrado en 0,5. Fuente: decision_automl_prep.json / panel mm25.]
\begin{figure}[htbp]
\centering
\fbox{\parbox{0.92\textwidth}{\centering\vspace{1.0cm}\textit{Heatmap de accuracy (activo $\times$ estrategia) y cuota SHAP por bloque, exportar desde el notebook §4.4.}\vspace{1.0cm}}}
\caption{Izquierda: accuracy por activo y estrategia centrada en el azar. Derecha: importancia media por bloque de características, con STRATA por delante del agente.}
\label{fig:mp-panel-heatmap}
\end{figure}

\section{Patrón entre activos: ¿qué estrategia conviene a cada tipo?}
\label{sec:mp-clustering}

No todos los activos son iguales: cambian el \emph{leverage effect}, la volatilidad y el sesgo del agente. Agrupamos los quince por su naturaleza, con características de mercado y no de rendimiento para no contaminar la lectura, mediante cuatro métodos adecuados a una muestra pequeña: KMeans, Ward, mezcla gaussiana y agrupamiento espectral. La calidad de la partición favorece tres grupos: la silueta es máxima en $k=3$ para los tres métodos concordantes ($0{,}485$ en KMeans, Ward y la mezcla gaussiana) y el criterio BIC de la mezcla gaussiana cae con holgura al pasar de dos a tres grupos. El método espectral discrepa, y lo recogemos como tal.

Los tres grupos admiten lectura económica. El primero reúne los índices y sectoriales de \emph{leverage} fuerte (SPY, QQQ, DIA, IWM, XLE, XLF, XLK, BAC), con correlación retorno-volatilidad media de $-0{,}09$. El segundo agrupa los volátiles de apalancamiento débil (NVDA, TSLA, ROKU, MARA). El tercero recoge tres activos atípicos (MSTR, SMCI, UNG). Cruzando los grupos con la Tabla~\ref{tab:mp-panel} emerge una hipótesis: en los índices de \emph{leverage} fuerte el valor está en el canal de régimen, que rescata al agente y disciplina su riesgo aunque no bata a las triviales; en los volátiles y atípicos el valor está en el meta-learner y en AutoML, que sí superan a las triviales. La conclusión es exploratoria, a una muestra de quince activos, y la presentamos como material para futuras líneas (Figura~\ref{fig:mp-cluster}).

% [Figura por exportar — notebook §4.5, celda 22: proyección PCA 2D de la naturaleza de los activos,
%  coloreada por grupo (k=3). Fuente: strategy_clustering15.json.]
\begin{figure}[htbp]
\centering
\fbox{\parbox{0.92\textwidth}{\centering\vspace{1.0cm}\textit{Proyección PCA 2D de la naturaleza de los activos por grupo ($k=3$), exportar desde el notebook §4.5.}\vspace{1.0cm}}}
\caption{Agrupamiento de los quince activos por naturaleza de mercado (proyección a dos componentes principales). Los índices de \emph{leverage} fuerte, los volátiles y los atípicos se separan.}
\label{fig:mp-cluster}
\end{figure}

\section{Robustez y honestidad}
\label{sec:mp-robustez}

Conviene poner la señal de STRATA a prueba contra una alternativa exigente. Sobre SPY, partiendo de un modelo con solo características de momento, añadir las siete señales de STRATA sube la accuracy de $0{,}521$ a $0{,}582$, una mejora de $+0{,}06$ que se repite en los cinco bloques de semillas y resulta significativa en el McNemar de tres. La señal de STRATA aporta información que el momento no captura.

El cuadro cambia en SMCI, una acción de apalancamiento débil. Ahí el régimen de Crisis tiene media positiva, el HMM pierde su contenido direccional y el margen de rescate se estrecha. El M10 desplegable acierta el $55{,}2\%$ y gana a todas las referencias, pero de forma nominal: el binomial frente a la clase mayoritaria da $p = 0{,}14$. El límite estaba previsto en el diseño, ya que la supervisión rescata con fuerza cuando el agente discrepa de un régimen que acierta, y SMCI cumple poco esa condición.

Queda claro qué sobrevive a un test y qué no. El rescate del agente en accuracy es significativo (McNemar de M10 y AutoML frente a M5 sobre SPY, $p = 0{,}007$ y $p = 0{,}0002$). El rescate en riesgo lo es en el agregado del panel (Sharpe y máxima caída de M8 frente a M5, con intervalos que excluyen el cero). Batir a las triviales en accuracy no alcanza significancia, sin potencia a doscientos cincuenta días. La señal informativa reside en STRATA, con una cuota de importancia media del $66\%$. Todo se mide con \texttt{signal\_lag}${}=1$, embargo de un día, modelos fijados \emph{ex ante} y contrastes con cita, y el notebook cierra con un auto-test que cruza cada cifra de cabecera con su fichero de origen.

\section{Conclusiones del marco práctico}
\label{sec:mp-conclusiones}

Cada afirmación del capítulo queda sostenida por la evidencia que admite. El agente solo pierde y acierta menos de la mitad (SPY: accuracy $0{,}366$, Sharpe muy negativo, sign test que rechaza el azar). STRATA lo rescata, y el rescate se prueba: en accuracy con el McNemar frente al agente, y en riesgo con el bootstrap agregado del panel, donde la regla M8 mejora el Sharpe en $+0{,}66$ y la máxima caída en $+0{,}24$ con intervalos que excluyen el cero. Un aprendiz potente reencuentra la señal de STRATA y se apoya en ella, con una cuota de importancia media de dos tercios, y sobre un modelo de momento esas señales añaden acierto. El agrupamiento por naturaleza sugiere un patrón entre el tipo de activo y la estrategia que le conviene, que dejamos como hipótesis exploratoria. El límite se dice con la misma claridad: ninguna estrategia bate a las triviales en accuracy con significancia a esta muestra, y donde el apalancamiento es débil el margen de rescate se estrecha.

La aportación es un protocolo de supervisión estadística interpretable que recupera a un agente perdedor en acierto y en riesgo, que un meta-learner potente confirma al apoyarse en sus señales, y que delimita con honestidad dónde funciona. La significancia plena de la ventaja en accuracy y la generalización a otros agentes y activos quedan como las líneas inmediatas del Capítulo~\ref{cap:conclusiones}.
